\documentclass{article}
\usepackage{amsmath, amssymb, amsthm}
\usepackage{hyperref}
\usepackage{geometry}
\geometry{margin=1in}

\title{Erd\H{o}s Problem \#789}
\date{Last updated: 2025-08-31}
\author{Source: erdosproblems.com}

\begin{document}
\maketitle

\noindent\textbf{Status:} OPEN \quad \textbf{No prize}

\noindent\textbf{Tags:} additive combinatorics

\medskip

\noindent\textbf{Problem Statement:}

Let $h(n)$ be maximal such that if $A\subseteq \mathbb{Z}$ with $\lvert A\rvert=n$ then there is $B\subseteq A$ with $\lvert B\rvert \geq h(n)$ such that if $a_1+\cdots+a_r=b_1+\cdots+b_s$ with $a_i,b_i\in B$ then $r=s$. Estimate $h(n)$.




Erd\H{o}s \cite{Er62c} proved $h(n) \ll n^{5/6}$. Straus \cite{St66} proved $h(n) \ll n^{1/2}$. Erd\H{o}s noted the bound $h(n)\gg n^{1/3}$, taking\[B=\{ a: \{ \alpha a\} \in n^{-1/3}+\tfrac{1}{2} (-n^{-2/3},n^{-2/3})\}\]for a random $\alpha\in [0,1]$. \cite{Er62c} and Choi \cite{Ch74b} improved this to $h(n) \gg (n\log n)^{1/3}$. See also [186] and [874] .




References

[Ch74b] Choi, S. L. G., On an extremal problem in number theory . J. Number Theory (1974), 105--111.

[Er62c] Erd\H{o}s, P\'{a}l, Some remarks on number theory. {III} . Mat. Lapok (1962), 28--38.

[St66] Straus, E. G., On a problem in combinatorial number theory . J. Math. Sci. (1966), 77--80.

\noindent\textbf{Related OEIS sequences:} possible


\bigskip
\noindent\small{Source: \url{https://www.erdosproblems.com/789}}
\end{document}
